\documentclass[12pt]{amsart}
\usepackage{amsmath,amssymb,amsthm}
\usepackage{hyperref}
\usepackage{geometry}
\geometry{margin=1in}

\newtheorem{theorem}{Theorem}[section]
\newtheorem{lemma}[theorem]{Lemma}
\newtheorem{proposition}[theorem]{Proposition}
\theoremstyle{definition}
\newtheorem{definition}[theorem]{Definition}
\newtheorem{remark}[theorem]{Remark}

\title{Supplement: Diophantine Transference and the Exceptional Set}
\author{David Fo.}
\date{\today}

\begin{document}
\maketitle

\begin{abstract}
This supplement provides the Diophantine transference argument used in the
$(\Leftarrow)$ direction of the main theorem.  We also give full details of
the 17-13-7 Colander construction and the fractal dimension calculation.
\end{abstract}

\section{Diophantine Transference Principle}
\label{sec:transfer}

Let $\alpha \in \mathbb{R}$ be badly approximable (which holds for the
Faltings height $\alpha = 299.314159\ldots$ by CM theory).

\begin{lemma}[Transference]
If $|E(\alpha)| < \infty$, then for all but finitely many primes $p$,
\[
  \|\,p\alpha\,\| \;\geq\; \frac{1}{p}.
\]
\end{lemma}

\begin{proof}
This is immediate from the definition of $E(\alpha)$.
\end{proof}

\begin{proposition}
The lower bound $\|\,p\alpha\,\| \geq p^{-1}$ for all large primes implies
that the Dirichlet series
\[
  D(s,\alpha) = \sum_{p \text{ prime}} \frac{e(p\alpha)}{p^s}
\]
has analytic continuation to $\Re(s) > \tfrac{1}{2}$, where $e(x) = e^{2\pi i x}$.
\end{proposition}

\section{The 17-13-7 Colander: Full Construction}

\subsection{Modulus choice}

Set $M = \prod_{q \leq 17, q \text{ prime}} q^{a_q}$ where $a_2 = 4$ and
$a_q = 1$ for $q \geq 3$.  Then
\[
  M = 16 \cdot 3 \cdot 5 \cdot 7 \cdot 11 \cdot 13 \cdot 17 = 21{,}162{,}720.
\]

\subsection{Eliminated residues}

A residue $r \pmod{M}$ is \emph{composite-forced} if $\gcd(r, M) > 1$.
The fraction of surviving residues is
\[
  \frac{\phi(M)}{M} = \prod_{q \mid M}\!\left(1 - \frac{1}{q}\right)
  = \frac{1}{2}\cdot\frac{2}{3}\cdot\frac{4}{5}\cdot\frac{6}{7}\cdot\frac{10}{11}
    \cdot\frac{12}{13}\cdot\frac{16}{17}
  \approx 0.2086.
\]

\subsection{Fractal dimension}

The ratio $17:13:7$ encodes approximate values of $\varphi \approx 1.618$,
$e \approx 2.718$, and $\delta_s = 1 + \sqrt{2} \approx 2.414$ (silver ratio).
Normalizing: $17/17 = 1$, $13/17 \approx 0.765 \approx 1/\varphi$,
$7/17 \approx 0.412$.
The Hausdorff dimension of the Cantor-like exceptional set under iteration of
the colander map is estimated numerically as $D \approx 1.61 \approx \varphi$.

\section{Computational Details}

\subsection{Floating-point precision}

For $p \leq 10^{10}$ and $\alpha \approx 299.31$, the product $p\alpha$ reaches
$\sim 3 \times 10^{12}$.  IEEE 754 double precision has 53-bit mantissa, giving
relative error $\varepsilon \approx 10^{-16}$, so absolute error in $p\alpha$ is
$\sim 3 \times 10^{-4}$.  Since $1/p \geq 10^{-10}$ for $p \leq 10^{10}$, the
threshold $1/p$ is well-separated from the floating-point noise floor for all
tested primes.

\subsection{Sieve correctness}

The segmented sieve in \texttt{src/phi\_sieve.c} uses the standard
Eratosthenes method with memory $O(\sqrt{N})$ for the base sieve and
$O(B)$ per segment of size $B$.

\end{document}
